\section{Reading the active-space verification}
\label{sec:manifold-tutorial}

At the robust level the driver verifies the selected active space by
the entanglement entropy of the converged wavefunction and refuses the
run when the space is single-reference.  This chapter shows what the
verification prints, how to read it, and how to inspect the entropy of
a spin multiplet directly.  The manual chapter
\emph{Manifold-merged entanglement entropy} carries the derivation; this
chapter is about using it.

\subsection{A closed-shell case is refused, and that is the point}

The water molecule with a CAS(2,2) over the two frontier orbitals is
weakly correlated at its equilibrium geometry.  Run it at the robust
level:
\begin{lstlisting}[language=Python]
from diffmrc.driver.kernel import compute_ground_state

out = compute_ground_state(
    'O 0 0 0; H 0 0.757 0.587; H 0 -0.757 0.587',
    charge=0, spin=0, active_space='4,5',
    params={'basis': '6-31g'}, level='robust')

data = out['algorithm_data']
print(data['active_space_verified'])
print(data['active_space_entropy']['s1'])
print(data['active_space_entropy']['verdict'])
\end{lstlisting}
The output is
\begin{lstlisting}
False
[0.004787230889527457, 0.004787230889527577]
single_reference
\end{lstlisting}
The largest entropy, $4.79\times10^{-3}$, lies far below the $0.14$
multireference criterion, so no active space is justified for this
system.  The robust level reports the run as unverified and refuses it.
This is not a failure of the calculation: the CASSCF energy is the same
as at the fast level, $E=-75.98495118733113$~Ha, to the last digit.  It
is the verification saying that the question the level promises to
answer does not have a positive answer here.

Use the fast level for such systems, or pass an active space
explicitly and read the energy without the verification.

\subsection{Why the verification needs the manifold}

A high-spin open shell is the case the verification exists for, and it
is also the case where a naive read of the entropy misleads.  Consider
the nitrogen atom in its $^{4}S$ ground state, three electrons in the
three $2p$ orbitals:
\begin{lstlisting}[language=Python]
import numpy as np
from diffmrc.core.nevpt2.impl import cas_dets_for2
from diffmrc.core.orblib import orbital_pairs
from diffmrc.driver.kernel import (casci_vectors_at,
                                   partition_from_occupations,
                                   scf_reference)
from diffmrc.analysis.ci import manifold_entropy_from_ci

mol, mf = scf_reference('N 0 0 0', 0, 3, '6-31g')
occ = np.asarray(mf.mo_occ)
C = np.asarray(mf.mo_coeff)
space = partition_from_occupations(occ, [2, 3, 4])

pairs = orbital_pairs(space.ncore, space.ncas, mol.nao_nr(),
                      include_cv=True)
vectors = casci_vectors_at(mol, C, space, np.zeros(len(pairs)),
                           nroots=4)          # 2S + 1

dets = cas_dets_for2(space.ncore, space.ncas, space.n_elec)
for k in range(vectors.shape[1]):
    s1 = manifold_entropy_from_ci(vectors[:, k], dets,
                                  ncore=space.ncore, nact=space.ncas)
    print(f'root {k}: {s1}')

merged = manifold_entropy_from_ci(vectors, dets,
                                  ncore=space.ncore, nact=space.ncas)
print(f'merged: {merged}')
\end{lstlisting}
The output is
\begin{lstlisting}
root 0: [0.660849 0.660849 0.660849]
root 1: [0.636514 0.636514 0.636514]
root 2: [2.235612e-25 1.548999e-24 5.149988e-25]
root 3: [0.168823 0.168823 0.168823]
merged: [0.693147 0.693147 0.693147]
\end{lstlisting}
Read one root at a time and the answer ranges from zero to $0.66$,
with nothing to choose between them.  Root 2 is the $M_S=+3/2$
determinant: all three $2p$ orbitals carry one $\alpha$ electron, so
each is in a definite Fock state and its four-state entropy is zero.
The other roots are the remaining projections of the same quartet.

The merged value is $\ln 2 = 0.693147$ for all three orbitals: one
electron of unknown spin per orbital, which is the physical content of
the degenerate open shell.  The gate uses this merged value, so the
space is certified as multireference.

\subsection{What the report contains}

\texttt{active\_space\_entropy} is a dictionary with five entries:

\begin{center}
\begin{tabular}{ll}
\toprule
key & meaning \\
\midrule
\texttt{s1} & merged single-orbital entropy of the active orbitals \\
\texttt{nroots} & number of components merged, $2S+1$ \\
\texttt{verdict} & the comparison against the selected space \\
\texttt{plateau\_heights} & detected plateaus of the merged spectrum \\
\texttt{evidence} & human-readable sentences behind the verdict \\
\bottomrule
\end{tabular}
\end{center}

The verdict is \texttt{agrees}, \texttt{candidate\_too\_small},
\texttt{candidate\_too\_large}, \texttt{no\_plateau},
\texttt{single\_reference}, or \texttt{failed}.  Only
\texttt{single\_reference} vetoes the run; the size comparisons are
recorded because the element rules deliberately add ligand orbitals
whose entropy lies below the line.

\subsection{Two limits to keep in mind}

The spectrum covers the selected space only.  An orbital that the
selection left out contributes zero by construction, so the gate
cannot see a truncated frontier; that remains the job of the selection
rules.

The merge assumes that the lowest $2S+1$ roots are the components of the
target manifold, which holds when they are degenerate.  Print the root
energies when the assumption is in doubt.

Finally, the end-to-end driver on the nitrogen quartet takes a
different route to the curvature.  The converged point sits within
$2\times10^{-7}$ of the symmetric $\theta=0$, where the exact
degeneracy of the quartet makes every entry of the analytic orbital
Hessian non-finite.  The energy and the analytic gradient are finite
there, so the driver takes the curvature from a symmetric difference of
the gradient at a step of $10^{-5}$ (capped at 300 rotation parameters;
a larger space is refused rather than approximated) and records
\texttt{curvature\_method = 'gradient-difference'} in the result.  The
listing above calls \texttt{casci\_vectors\_at} directly and is
unaffected by any of this.
